% Agent 22: Pages 427--430 (PDF pp.42--44)
% Sections 5.5 (cont.) and 5.6: Underlying Assumptions & Equations of Radial Structure

%%% ============================================================
%%% Continuation of Section 5.5 — Relativistic Model for Disk:
%%% Underlying Assumptions
%%% ============================================================

In calculating the disk structure, we shall make the following assumptions
and idealizations:
\begin{enumerate}
\item[(i)] The central plane of the disk coincides with the equatorial plane of the
black hole.
\item[(ii)] The companion star in the binary system has negligible gravitational
influence on the disk structure (assumption valid in inner part of disk;
not at outer edges).
\item[(iii)] The disk is thin; i.e., its proper thickness $2h$, as measured in the orbiting
frame of the disk's matter (same as proper thickness measured by any
other observer who moves in the equatorial plane), satisfies
\begin{equation}
h(r) \ll r. \tag{5.5.1}
\end{equation}
\item[(iv)] The disk is in a quasisteady state. This assumption can be made precise
in terms of an averaging process. Pick a particular height $z$ ($|z| \leq h \ll r$)
above the central plane of the disk, and pick a particular quantity $\Psi$
(e.g.\ $\Psi = {}$density or temperature of 4-velocity of gas) to be measured at
height~$z$. Average $\Psi$ over all $\phi$ (Lie dragging it along $\partial/\partial\phi$ while averaging,
if it is a vector or tensor), average over a proper radial distance of order
$2h$, and average over the time interval required for gas to move inward
a distance $2h$:
\begin{equation}
\langle \Psi \rangle \equiv \frac{1}{2\pi \Delta r \Delta t}
\int_{r-\epsilon r/2}^{r+\epsilon r/2}\int_{t-\epsilon t/2}^{t+\epsilon t/2}
\left(\frac{\mathscr{g}^{1/2}}{|\mathscr{g}\!|^{1/2}/\mathscr{B}\mathscr{g}^{1/2}}\right)
\;\Psi\, d\phi\, dt\, dr; \qquad \Delta r = 2h,
\tag{5.5.2}
\end{equation}
($\mathscr{g}^2$ will be defined below).  The assumption of a quasisteady state means
that $\langle \Psi \rangle$ is time-independent---i.e.,
\begin{equation}
\partial \langle \Psi \rangle / \partial t = 0 \quad \text{if } \langle \Psi \rangle \text{ is a scalar or tensor,}
\tag{5.5.3}
\end{equation}
\begin{equation}
\mathscr{L}_{\partial/\partial t}\langle \Psi \rangle = 0 \quad \text{if } \langle \Psi \rangle \text{ is a vector or tensor,}
\tag{5.5.3}
\end{equation}
where $\mathscr{L}_{\partial/\partial t}$ is the Lie derivative along the Killing vector $\partial/\partial t$.  The
disk structure may be very far from steady state on scales $\lesssim h$; for
example, it may have turbulence, flux reconnection, and flares on such
scales.  Such local violence will not invalidate the analysis that follows.
\item[(v)] When viewed ``macroscopically'' (turbulence removed by the above
averaging process), the gas of the disk moves (very nearly) in direct,
circular, geodesic orbits---i.e., with 4-velocity ($\mathbf{u} \simeq \mathbf{e}_6$). Superimposed
on this orbital motion is a very small radial flow, produced by viscous
stresses, and an even smaller flow in the vertical direction, as required
by the variation of disk thickness with radius.
\begin{equation}
\langle u \rangle = \frac{\mathbf{e}_6}{[1 - (\tilde{v})^2 - (\tilde{v}^{\hat{r}})^2]^{1/2}}
+ \tilde{v}^{\hat{\phi}}\,\mathbf{e}_{\hat{\phi}}
+ \tilde{v}^{\hat{r}}\,\mathbf{e}_{\hat{r}},
\tag{5.5.4}
\end{equation}
\end{enumerate}

\noindent
where $|\tilde{v}^{\hat{r}}(r,z)| \ll |\tilde{v}_{\hat{\phi}}| \simeq (M/r)^{1/2}$.

By assuming that the orbital motion of the gas is very nearly geodesic,
we are automatically requiring that the gravitational pull of the hole
dominates over radial pressure gradients and over stresses.  In
order of magnitude requiring that the gravitational pull of the hole (``gravitational
acceleration'') is the gradient of the gravitational binding energy,
$(1-\tilde{E})$; and the acceleration due to pressure gradients and shears is
$\sim (r^{\hat{k}}/\rho_0) \sim (r^{\hat{k}}\!/\rho_0)_r$.  Thus, we are automatically demanding that the
stress tensor in the fluid's rest frame (${}={}$orbiting orthonormal frame)
satisfy
\begin{equation}
\frac{r^{\hat{i}\hat{k}}(r,z)}{\rho_0(r,z)} \ll 1 - \tilde{E}(r) \simeq \frac{1}{2}\,\frac{M}{r}\,.
\tag{5.5.6}
\end{equation}

(Here $\simeq$ means ``equals in the Newtonian limit''.)

Because the specific internal energy $\Pi$ is approximately equal to
$r^{\hat{k}}\!/\rho_0$, the assumption of nearly geodesic orbits implies the condition
of ``negligible specific heat'':
\begin{equation}
\Pi(r,z) \ll 1 - \tilde{E}(r) \simeq \frac{1}{2}\,\frac{M}{r}\,.
\tag{5.5.7}
\end{equation}

In words---the density of internal energy, $\rho_0\Pi$ (including thermal
energy of gas, energy of turbulent motions, and magnetic-field energy)
is much smaller than the density of gravitational binding energy
$\rho_0(1-\tilde{E})$.

Notice that conditions (5.5.6) and (5.5.7) imply
\begin{equation}
\Pi \ll 1, \qquad T\mathscr{g}/\rho_0 \ll 1.
\tag{5.5.8}
\end{equation}

at all $r$ and $z$---even near the black hole.  This means that, although
general relativistic effects (spacetime curvature; deviations from
Newtonian gravity) are very important near the hole, one can everywhere
ignore special relativistic corrections to the local thermodynamic,
hydrodynamic, and radiative properties of the gas.  Contrast this
situation with the interior of a neutron star and the early stages of the
Universe, where not only is gravity relativistic but so is the matter
($\Pi \geq 1$, $p/\rho_0 \geq 1$, $T\mathscr{g}/\rho_0 \geq 1$).

Turn now from the underlying assumptions of the model to the reference
frames to be used in the calculations.  Three reference frames will be needed:
\begin{enumerate}
\item[(i)] The Boyer--Lindquist coordinate frame $\{(t,r,\theta,\phi)\}$ inside and near disk.
$(t,r,\phi)$ inside and near disk;
see eqs.\ (5.4.2a).
\item[(ii)] The orbiting

orthonormal frame $[\mathbf{e}_{\hat{0}}, \mathbf{e}_{\hat{r}}, \mathbf{e}_{\hat{z}}]$; eqs.\ (5.4.5).
\item[(iii)] The mean local rest-frame of the gas
\end{enumerate}

\begin{equation}
\bar{\mathbf{e}}_{\hat{0}} = \langle \mathbf{u} \rangle \quad \text{as given by eq.\ (5.5.4),}
\tag{5.5.9}
\end{equation}
\begin{align}
\bar{\mathbf{e}}_{\hat{z}} &= \mathbf{e}_{\hat{z}}, \notag\\
\bar{\mathbf{e}}_{\hat{\phi}} &= (\mathbf{e}_{\hat{\phi}} + \tilde{v}^{\hat{\phi}}\mathbf{e}_{\hat{0}})/
|\mathbf{e}_{\hat{\phi}} + \tilde{v}^{\hat{\phi}}\mathbf{e}_{\hat{0}}|, \notag\\
\bar{\mathbf{e}}_{\hat{r}} &= (\mathbf{e}_{\hat{r}} + \tilde{v}^{\hat{r}}\,\mathbf{e}_{\hat{0}}
- \tilde{v}^{\hat{r}}\,\tilde{v}^{\hat{\phi}}\,\mathbf{e}_{\hat{\phi}})
/|\mathbf{e}_{\hat{r}} + \tilde{v}^{\hat{r}}\,\mathbf{e}_{\hat{0}} - \tilde{v}^{\hat{r}}\,\tilde{v}^{\hat{\phi}}\,\mathbf{e}_{\hat{\phi}}|.
\tag{5.5.9}
\end{align}

The mean local rest frame is nearly identical to the orbiting orthonormal frame.


%%% ============================================================
%%% Section 5.6 — Equations of Radial Structure
%%% ============================================================

\subsection{Equations of Radial Structure}
\label{sec:5.6}

The relativistic equations of radial structure for our disk will be expressed in
terms of the following parameters: (i)~The steady-state accretion rate $\dot{M}_0$;
(ii)~The surface density of the disk
\begin{equation}
\Sigma \equiv \int_{-h}^{+h} \langle \rho_0 \rangle\, dz,
\tag{5.6.1a}
\end{equation}
where $\rho_0$ is the density of rest mass as measured in the rest frame of the gas.
(iii)~The integrated shear stress
\begin{equation}
W \equiv \int_{-h}^{+h} \langle r\hat{\phi}\hat{z} \rangle\, dz,
\tag{5.6.1b}
\end{equation}
where $T_{\hat{\phi}\hat{z}}$ is the component of the stress-energy tensor on the orbiting ortho\-normal basis vectors $\mathbf{e}_{\hat{z}}$ and $\mathbf{e}_{\hat{\phi}}$.  (iv)~The mass-averaged radial velocity of the gas
\begin{equation}
\bar{v}^{\hat{r}} \equiv (1/\Sigma) \int_{-h}^{+h} \langle v^{\hat{r}} \rho_0 \rangle\, dz.
\tag{5.6.1c}
\end{equation}
(v)~The flux of radiant energy off the upper face of the disk (equal also to flux
off lower face)
\begin{equation}
F = \langle r^{\hat{0}\hat{z}}(z = h) \rangle
= \langle -r^{\hat{0}\hat{z}}(z = -h) \rangle.
\tag{5.6.1d}
\end{equation}

\noindent (vi)~The mass $M$ and specific angular momentum $a$ of the hole.  (vii)~The radial
coordinate~$r$.

The laws governing the radial structure are conservation of rest mass,
conservation of angular momentum, and conservation of energy.

\bigskip
\noindent\textit{Conservation of rest mass}

\medskip
The amount of rest mass that flows inward across of cylinder a radius $r$ during
coordinate time $\Delta t$, when averaged by the method of equation (5.5.2), must
equal $\dot{M}_0 \Delta t$ (conservation of rest mass).  The mass transferred can be written
as the flux integral
\begin{equation}
\dot{M}_0 \Delta t = \int_{\mathscr{S}} \langle \rho_0 \mathbf{u} \rangle \cdot d^3\Sigma
= (-2\pi r \mathscr{g}^{1/2} \Delta t)
\int_{-h}^{+h} \langle \rho_0 \tilde{v}^{\hat{r}} \rangle\, dz,
\tag{5.6.2}
\end{equation}
where $\mathscr{S}$ is the 3-surface $\{0 \leq \phi \leq 2\pi,\; -h \leq z \leq h,\; 0 \leq t \leq \Delta t\}$.
Rewritten in terms of the mass-averaged radial velocity [eq.\ (5.6.1c)], this equation becomes
\begin{equation}
\dot{M}_0 = -2\pi r\, \mathscr{g}^{1/2}\, \Sigma\, \bar{v}^{\hat{r}}
\quad \text{(constant, independent of $r$ and $t$)}.
\tag{5.6.2}
\end{equation}

\bigskip
\noindent\textit{Conservation of angular momentum}

\medskip
The law of conservation of angular momentum can be written in the form
\begin{equation}
\mathbf{J} = \mathbf{T} \cdot (\partial/\partial\phi),
\qquad \nabla \cdot \mathbf{J} = 0,
\tag{5.6.3}
\end{equation}
where $\partial/\partial\phi$ is the Killing vector associated with rotation about the symmetry
axis.  Without any loss of generality, we can write the stress-energy tensor $\mathbf{T}$ in
the form
\begin{equation}
\mathbf{T} = \rho_0 (1+\Pi)\,\mathbf{u} \otimes \mathbf{u}
+ \mathbf{t} \cdot \mathbf{u} \otimes \mathbf{q}
+ \mathbf{q} \otimes \mathbf{u} + \mathbf{t},
\tag{5.6.4a}
\end{equation}
where $\Pi$ is specific internal energy, $\mathbf{t}$ is the stress tensor as measured in the local
rest frame of the baryons, $\mathbf{q}$ is the energy flux relative to the local rest frame
(see \S2.5), and
\begin{equation}
\mathbf{u} \cdot \mathbf{t} = \mathbf{u} \cdot \mathbf{q} = 0.
\tag{5.6.4b}
\end{equation}

The corresponding density of angular momentum is
\begin{equation}
p^\phi = \rho_0 u_\phi + q^\alpha u_{a\phi} + q_\alpha u^\alpha u_\phi
\tag{5.6.5}
\end{equation}
where we have dropped the momentum $\rho_0\Pi u_\phi + u^\alpha \mathbf{q} \cdot \nabla u_\phi + \mathbf{q} \cdot (q_\alpha u^\alpha)$
internal energy because $\Pi \ll 1$.  The law of angular-momentum conservation thus
reads
\begin{equation}
\nabla \cdot \mathbf{J} = 0 = \rho_0 (du_\phi/d\tau)_{;r} + r^{-1}(r^2 \mathscr{g}^{-1/2} \tilde{\Sigma} \tilde{v}^{\hat{r}} \tilde{L}_{,r})
+ r^{-1}(r^2 \mathscr{g}^{-1/2} 2\tilde{W})_{,r} + 2\tilde{L}F = 0.
\end{equation}

Inserting this and the relation
\begin{equation}
\langle du_\phi/d\tau \rangle = (u_\phi)_{,r}\, \mathscr{g}^{1/2} \bar{v}^{\hat{r}}
\tag{5.6.5}
\end{equation}
into the conservation law, obtain
\begin{equation}
\mathscr{g}^{1/2}\,\Sigma\,\bar{v}^{\hat{r}}\,\tilde{L}_{,r}
+ r^{-1}(r^2\mathscr{g}^{-1/2}2\mathscr{g}W)_{,r} + 2\tilde{L}F = 0.
\tag{5.6.5}
\end{equation}

The first term is the rate of increase of angular momentum in the gas; the second
is the rate at which shear stresses carry off angular momentum; the third is the
rate at which photons carry off angular momentum.  When combined with the
law of mass conservation (5.6.2), this law of angular momentum conservation
takes on the simpler form
\begin{equation}
-\dot{M}_0 \tilde{L}/2\pi + r^2 \mathscr{g}^{-1/2} 2\mathscr{g} W)_{,r}
+ 2\tilde{L}F = 0.
\tag{5.6.6}
\end{equation}

\bigskip
\noindent\textit{Conservation of energy}

\medskip
Turn now to the law of energy conservation
\begin{equation}
\mathbf{u} \cdot (\nabla \cdot \mathbf{T}) = 0.
\tag{5.6.7}
\end{equation}

Rewrite this law in the form
\begin{equation}
\nabla \cdot (\mathbf{u} \cdot \mathbf{T}) - u_{\alpha;\beta}\, r^{\alpha\beta} = 0.
\tag{5.6.8}
\end{equation}

Use the general stress-energy tensor (5.6.4) to write
\begin{equation}
\mathbf{u} \cdot \mathbf{T} = -\rho_0(1+\Pi)\,\mathbf{u} - \mathbf{q}.
\tag{5.6.9}
\end{equation}

and combine this with the law of rest-mass conservation $\nabla \cdot (\rho_0 \mathbf{u}) = 0$ to obtain
\begin{equation}
\nabla \cdot (\mathbf{u} \cdot \mathbf{T}) = -\rho_0\, d\Pi/d\tau
- \nabla \cdot \mathbf{q}.
\end{equation}

Decompose $u_{\alpha;\beta}$ into its irreducible tensorial parts
\begin{equation}
u_{\alpha;\beta} = \sigma_{\alpha\beta} + \tfrac{1}{3}\theta\, g_{\alpha\beta} + a_\alpha u_\beta
\end{equation}
[eq.\ (2.5.17)], and contract it into the stress-energy tensor (5.6.4a) to obtain
\begin{equation}
u_{\alpha;\beta}\, r^{\alpha\beta} = -\sigma_{\alpha\beta}\,r^{\alpha\beta}
+ \tfrac{1}{3}\theta\, r^\alpha{}_\alpha + \mathbf{a} \cdot \mathbf{q}.
\tag{5.6.10}
\end{equation}

Then use relations (5.6.9) and (5.6.10) to rewrite the law of energy conservation
(5.6.8) in the form
\begin{equation}
\rho_0\, d\Pi/d\tau + \nabla \cdot \mathbf{q}
= -\sigma_{\alpha\beta}\,r^{\alpha\beta}
+ \tfrac{1}{3}\theta\, r^\alpha{}_\alpha
+ \mathbf{a} \cdot \mathbf{q}.
\tag{5.6.11}
\end{equation}

The various terms in this equation have simple interpretations.  The left-hand side
represents the fate of the energy being generated locally in the gas: $\rho_0\, d\Pi/d\tau$ is
the energy being generated locally in the gas; $\nabla \cdot \mathbf{q}$ is the energy being transported out
of the region of generation.  The right-hand side represents the rate at which
energy is generated: $\sigma_{\alpha\beta}\,r^{\alpha\beta}$ is the energy being fed in by ``frictional''
(viscous) heating; $-\tfrac{1}{3}\theta$ is the energy being fed in by compression.  The remain-
